Small language models performing multi-step knowledge queries must continue from histories that depart from successful demonstrations. Extra executable actions can change intermediate-result identifiers and invalidate direct reuse of a reference suffix. We present dependency-consistent recovery supervision: a fixed initial model supplies unsuccessful training histories, reference dependencies are mapped to current result handles, and deterministic replay validates the continuation. Training retains the actual divergence as context and supervises subsequent assistant actions. Using Qwen2.5-3B-Instruct with low-rank adaptation, we compare recovery and ordinary continuation training under matched initialization, successful material, question order, and supervised-token exposure. On 2,000 project-held-out KQA Pro questions, mean accuracy across two continuation seeds increases from 78.75\% to 83.50\%, a gain of 4.75 percentage points with a question-paired 95\% bootstrap interval of $[3.53,5.98]$. Cumulative input-plus-generated tokens decrease by 8.84\% and 4.17\% relative to the matched stepwise controls. An offline diagnostic on 414 comparable historical training pairs finds that all suffix calls execute in 65 direct-splice cases and all 414 remapped cases. These results support reference-program-based recovery supervision for small models in deterministic, read-only knowledge-query environments.
